\subsection{ARCH-136-DEGREE-TWO-FREE: a free version of the concentration construction}

\paragraph{Archive locators.}
The remaining vertical condition is on cote 136, physical p.~117,
footer 116, \texttt{ARCHIVE\_136\_scan117.tex}; horizontal uniformity
is on physical p.~119, footer 118, \texttt{ARCHIVE\_136\_scan119.tex}.
The common-source question is physical p.~121, footer 120,
\texttt{ARCHIVE\_136\_scan121.tex}.

\paragraph{Statement.}
The source in ARCH-136-DEGREE-TWO can be chosen with every
$A_n^{w,q}$ free as an abelian group, and with
\[
             H^0(A_n^{w,*})=S^w,\qquad
             H^1(A_n^{w,*})=0                            \tag{136.DTF1}
\]
in every weight, rather than only an eventual range.  Its construction
uses the same fixed horizontal and higher-degree cutoffs, but omits
the relative $H^1$ repair entirely.  It still leaves arbitrary integer
torsion in degree two; it does not restore a finite-prime bound there.

\paragraph{Proof.}
Use the same target replacements with simplex zero exactly $S$,
and begin with the constant $S$.  Alternate only the horizontal
free-cell step and the reduced-representable Tate-pair step $V$
for degrees $q\geq3$.  The proof of ARCH-136-DEGREE-TWO already
establishes that these steps preserve simplex zero, all target maps,
weight zero and exact rational vertical cohomology.  The class-killing
argument for horizontal rows and vertical degrees $q\geq3$ therefore
applies without change.  It remains to check the stronger low-degree
assertion and freeness.

At a fixed simplex, a horizontal attachment is a free extension by
vertical unit disks, using the degreewise free complement of its
boundary inclusion.  The integral block calculation of
ARCH-136-DISK-SUPPORT shows that every even disk contracts onto
constants; an odd disk, after its unit block is contracted, has
nonconstant blocks starting in degree at least three.  Tensoring
with any nonnegative complex cannot change $H^0$ or $H^1$ through
those remaining blocks.  The actual disk contractions handle the
other summands without a flatness assumption.  Thus every horizontal
step preserves (136.DTF1).

For a higher Tate-pair attachment with $q\geq4$, both new generator
degrees $q-2,q-1$ are at least two.  The extension leaves degrees
zero and one, and their differential images, unchanged, so again
$H^0,H^1$ are unchanged.
For $q=3$ write its generators as $t_\alpha$ in degree one and
$c_\alpha$ in degree two, with
\[
             dt_\alpha=N_\alpha c_\alpha-b_\alpha,
             \qquad dc_\alpha=z_\alpha,
             \qquad N_\alpha>0.
\]
The subscript runs through a free basis of the added module in
this simplex, and also allows any simultaneous collection of such
attachments.  Every new degree-one element has a unique expression
\[
                 a+\sum_\alpha f_\alpha t_\alpha,
                       \qquad a\in A^1,\ f_\alpha\in A^0,
\]
with finitely many nonzero coefficients.  In its differential,
the coefficient of the independent basis vector $c_\alpha$ is
$N_\alpha f_\alpha$.  The terms $(df_\alpha)t_\alpha$, the old
$d a$, and $-f_\alpha b_\alpha$ contain no new $c$ generator.
If $A^0$ is torsion-free, a cycle consequently has every
$f_\alpha=0$.  Its cycles in degree one are exactly the old cycles,
and degree zero has not changed.  Hence its degree-one boundaries
are also the old ones.  This proves preservation of (136.DTF1)
for the remaining case.

The torsion-free hypothesis just used holds throughout, and more
is true.  Initially every group of $S$ is free over $\mathbf Z$.
Each horizontal attachment and each higher attachment is, in a
fixed simplex and forgetting the differential, a free graded PD
extension by free abelian generator modules.  Its divided-power
and exterior monomial basis makes it a free extension of the old
underlying abelian group.  The inclusion of the old group is split
by setting the new generators to zero as a graded algebra map;
this splitting need not commute with the differential.  Therefore
all successor groups are free.  At a union stage their compatible
old bases together with the successive monomial complements give
a free basis of the union.  This proves degreewise freeness,
including freeness of $A^0$, at every stage and in the limit.

Finally filtered colimits preserve the low cohomology identities
just proved.  Coefficients remain the specified ones, not an
unidentified copy of $S$, because all maps preserve their inclusion.
The alternating class-killing argument yields the other conclusions
of ARCH-136-DEGREE-TWO with its fixed cutoffs.\hfill$\square$
